\section{The main term}\label{sec:main-term}

The filling theorem (\cref{thm:filling}) needs one bound that is sharp to a factor $1+O(b^{-4})$: the link second moment of the open elements of a fixer
(\cref{thm:main-term}). The terms over an open reduction lie in the big cell of the Bruhat decomposition, so their class sums are sums over double cosets
$Uw_p\sigma U$ (\cref{ss:big-cell}). At a finite set of places each double coset sum is a Kloosterman sum times an orbital integral at infinity
(\cref{ss:kloosterman}). Siegel's formula, used as an upper bound, turns the sums over rational points into integrals over spheres at the finite places
(\cref{ss:siegel-top}), and a sieve over all places, entry by entry, removes the sums over the entries of the heads at the cost of a factor near $1$
(\cref{ss:sieve-systems,ss:sieve}). The section opens with the count of bad types used in \cref{thm:spanning-inputs}(b).

\subsection{Bad types are rare}\label{ss:bad-types}
Let $\mathbb k$ be a field with $q$ elements and $1<r<q$. Give every polynomial of degree $r$ over $\mathbb k$ the weight $\omega_{q,r}=(1-r/q)/q^r$. The idempotent of a tuple $c=(c_w)$
of labels over fields $\mathbb k_w$ is $e_c=\bigotimes_we_{c_w}$, with $e_{c_w}=\sum_g(\Xi_{c_w}(1)/\abs{G_w})\overline{\Xi_{c_w}(g)}\,g$ in $\C[G_w]$, $G_w=\GL_r(\mathbb k_w)$ (\cref{ss:B-green}).
\lean{New.IsAllowedPolynomial, New.IsPrimitivePolynomial, New.labelIdem, New.tupleIdem}

\begin{proposition}[Bad types]\label{prop:bad-types}
Let $M\ge1$ and $1<r<q$.
\begin{enumerate}[label=\textup{(\alph*)}]
\item At most $(1+M+r)q^{r-1}$ monic polynomials of degree $r$ over $\mathbb k$ are not allowed for $M$, and at least $\varphi(q^r-1)/r$ are primitive.
\item If $(1+M+2r)/q\le\varepsilon$, the allowed polynomials of degree $r$ weigh at least $1-\varepsilon$; the polynomials that are not primitive weigh at most
$\exp\bigl(-1/(2r(2r\log q+1))\bigr)$.
\item For labels $c_w$ of degree $r$ over fields $\mathbb k_w$ with $q$ elements, the coefficient of $e_c$ at the identity has real part at least $\prod_w\omega_{q,r}$.
\item Let there be $f$ fields $\mathbb k_w$, each with $q$ elements, and let $(1+M+2r)/q\le\varepsilon\le1$ and $2r(2r\log q+1)L\le f$. Then the coefficient at the identity of
$1-\sum_ce_c$, the sum over the good tuples $c$ for $M$, has real part at most $f\varepsilon+e^{-L}$.
\item For admissible data and $i=0,1$, the idempotent $e_{p_i}$, the sum of the $e_c$ over the good tuples of labels at the chosen places above $p_i$, is a central idempotent,
$e_0=e_{p_0}\otimes e_{p_1}$, and the coefficient of $1-e_{p_i}$ at the identity has real part at most $b^{-8}/2$.
\end{enumerate}
\end{proposition}
\lean{New.card_not_isAllowedPolynomial_le, New.totient_le_mul_card_isPrimitivePolynomial, New.one_sub_le_card_allowed_mul, New.card_not_primitive_mul_le, New.prod_le_re_tupleIdem_coeff_one, New.re_coeff_one_sub_sum_tupleIdem_le, New.Data.bad_types_rare, New.Data.re_coeff_one_sub_goodIdempotentAbove_le, New.Data.goodIdempotentAbove}

\begin{proof}
(a) A monic polynomial of degree $r$ that is not allowed either is not a label, that is, vanishes at $0$ or is not squarefree, or has a root $x$ in $\mathbb k$ with $x^M=1$. In the first and the last case it has
a root $x$ with $x=0$ or $x^M=1$; there are at most $1+M$ such $x$, and for each of them $q^{r-1}$ monic polynomials of degree $r$ vanish at $x$. A polynomial that is not squarefree is $P^2B$ with $P$, $B$
monic, $e=\deg P\ge1$ and $\deg B=r-2e<r$; for each $e\le r/2$ there are $q^e\cdot q^{r-2e}\le q^{r-1}$ such products. In all there are at most $(1+M+r)q^{r-1}$ polynomials that are not allowed. A primitive
polynomial of degree $r$ is the minimal polynomial of a generator of the cyclic group $\F_{q^r}^\times$: if a generator has minimal polynomial of degree $d$, then it lies in $\F_{q^d}$, so $q^r-1$ divides $q^d-1$ and
$d=r$. There are $\varphi(q^r-1)$ generators, and each primitive polynomial has $r$ of them as roots.
(b) The $q^r$ monic polynomials of degree $r$ weigh $q^r\omega_{q,r}=1-r/q$ together, and by (a) those that are not allowed weigh at most $(1+M+r)q^{r-1}/q^r=(1+M+r)/q$; so the allowed ones weigh at least
$1-r/q-(1+M+r)/q=1-(1+M+2r)/q\ge1-\varepsilon$. For the others let $P$ be the number of primitive polynomials, $m=q^r-1$ and $A=2r(2r\log q+1)$. Then
\[
  q^r\le2m\le2\varphi(m)(2\log m+1)\le2rP(2r\log q+1)=AP,
\]
by $m\le\varphi(m)(2\log m+1)$, by (a) and by $\log m\le r\log q$. The polynomials that are not primitive number $q^r-P$, and each weighs at most $q^{-r}$, so together they weigh at most $1-P/q^r\le1-1/A\le e^{-1/A}$.
(c) The coefficient of $e_c$ at $1$ is $\prod_w\Xi_{c_w}(1)^2/\abs{G_w}$, the coefficient of a tensor product at $1$ being the product of the coefficients. For a label $c$ with irreducible factors of degrees $d_i$, Green's degree formula gives
$\Xi_c(1)=\prod_{j=1}^r(q^j-1)/\prod_i(q^{d_i}-1)$ \cite{Green1955}, while $\abs{\GL_r(\mathbb k)}=q^{r(r-1)/2}\prod_{j=1}^r(q^j-1)$. Hence
\[
  \frac{\Xi_c(1)^2}{\abs{\GL_r(\mathbb k)}}=\frac{\prod_{j=1}^r(q^j-1)}{q^{r(r-1)/2}\prod_i(q^{d_i}-1)^2}\ge\frac{q^{r(r+1)/2}\bigl(1-\sum_{j=1}^{r}q^{-j}\bigr)}{q^{r(r-1)/2}q^{2r}}\ge\frac{1-r/q}{q^r},
\]
using $\sum_id_i=r$, $\prod_i(q^{d_i}-1)\le q^r$ and $\prod_j(1-q^{-j})\ge1-\sum_jq^{-j}\ge1-r/q$.
(d) The coefficient of $1$ at the identity is $1$, and by (c) that of $\sum_ce_c$ is at least $\omega^f$ times the number of good tuples, $\omega=\omega_{q,r}$. A tuple of allowed polynomials is good unless none of its
members is primitive. Weighted by $\omega^f$, the tuples of allowed polynomials count at least $(1-\varepsilon)^f\ge1-f\varepsilon$ by (b), and those without a primitive member at most $\exp(-1/A)^f=\exp(-f/A)\le e^{-L}$, with $A=2r(2r\log q+1)$ as in (b), since
$AL\le f$. Hence the coefficient of $1-\sum_ce_c$ has real part at most $1-(1-f\varepsilon-e^{-L})=f\varepsilon+e^{-L}$.
(e) The idempotents are central by \cref{thm:B-green}, and their tensor product is $e_0$ by the definition of the good families (\cref{prop:e0}). At the chosen places above $p_i$ the residue fields have $q=p_i$
elements and there are $f$ of them; with the order $q_2$ of the characters in place of $M$, the rank $r$ of $Y_0$ and $L=9\log b$, the size conditions of the data give $1<r<p_i$, $(1+q_2+2r)/p_i\le1$,
$2r(2r\log p_i+1)\cdot9\log b\le f$ and $f(1+q_2+2r)/p_i+b^{-9}\le b^{-8}/2$, and (d) applies with $\varepsilon=(1+q_2+2r)/p_i$.
\end{proof}

\subsection{Open elements, the big cell and orbit sums}\label{ss:big-cell}
For a tuple $p=(p_0,\dots,p_{\varrho-1})$ of units the \emph{Weyl element} $w_p$ sends $e_a$ to $p_ae^*_a$ and $e^*_a$ to $e_a/\bar p_a$ for $a<\varrho$, and fixes $X_0$ and $z$. The
\emph{floor group} $G^{(\varrho)}$ is the group of isometries that fix the $e_a$ and $e^*_a$, $a<\varrho$. A \emph{head} is a tuple $p$ followed by a last entry $t$. The \emph{entry set} $\mathcal T$ is
the finite set of quotients of two elements of the finite set of nonzero integers $\delta$ of $E$ with $\abs\delta_v^2\le b^{270n^2}$ at every infinite place. For
rational $\gamma$ in the floor group and a class pair $(c,c')$, the \emph{orbit sum} $\mathcal A^U(c,s,c';\gamma)$ is the sum of the amplitudes of $u\gamma u'^{-1}$ over the pairs $(u,u')$ of
rational points of $U$ for which $u\gamma u'^{-1}$ is a term over $(c,s,c')$.
\lean{New.Floors.FloorGroups.weylElem, New.Shape.Form.floorGroup, New.Filling.boundedNonzeroIntegers, New.Filling.finite_boundedNonzeroIntegers, New.Filling.finite_quotient_boundedNonzeroIntegers, New.Filling.entrySet, New.Filling.cosetAmplitude, New.Filling.translateByClassReps, New.Filling.NamesLink, New.padTuple}

\begin{lemma}[Open elements lie in the big cell]\label{lem:big-cell}
Let $\gamma$ be a term over $(c,s,c')$. For each $l$ the leading minor of order $l$ of the matrix $Z(\gamma)$ of frame products is an algebraic integer, and at every chosen place $w$
its reduction modulo $w$ is the same minor of the frame Gram matrix of $s$ at $w$. Hence if $s$ is open, no proper leading minor of $Z(\gamma)$ vanishes.
\end{lemma}
\lean{New.Filling.leadingMinor_frameGram_map, New.Filling.exists_leadingMinor_reduceAt_eq, New.Filling.IsOpenElement.leadingMinor_ne_zero}

\begin{proof}
$Z(\gamma)$ is the matrix of the products of the frame with $h_c^{-1}\gamma h_{c'}$ applied to the frame, up to the integral change of basis by $h_c$ and $h_{c'}$, and $h_c^{-1}\gamma h_{c'}\in K$
has integral entries in the lattice basis; so its minors are integral at every finite place, hence algebraic integers. Its reduction at $w$ is $s_w$, and taking minors commutes with the ring map to the residue
field. If $s$ is open, at some chosen place above $p_0$ no proper leading minor of the frame Gram matrix of $s$ vanishes; a minor whose reduction is nonzero is nonzero.
\end{proof}

\begin{lemma}[Weyl elements and the norm equation]\label{lem:weyl-norm}
Let $\sigma$ be rational in the floor group, $t=h(z,\sigma z)$ the last diagonal entry of $Z(\sigma)$ and $\omega$ the coordinates of $\sigma z$ in $X_0$.
\begin{enumerate}[label=\textup{(\alph*)}]
\item $N\omega^*\omega=t\bar t-N^2$; if $t\bar t\ne N^2$, then $\abs t_v>N$ at every infinite place.
\item $h_c^{-1}(w_p\sigma)h_{c'}=w_p(h_c^{-1}\sigma h_{c'})$, and $h_c^{-1}\sigma h_{c'}$ lies in the floor group; at $T$ the translate sends $z$ to $x-N^{-1}tz$ in the lattice basis, where $x$ is the part in $X_0$ of $\sigma z$
moved by the inverse of the representative of $c$.
\item At an infinite place $v$ the rescaled image of $w_p\sigma$ is $w_q$ times that of $\sigma$, with $q_a=v(p_a)$ times the squared scale of the frame vector $a$.
\item The link of $uw_p\sigma u'^{-1}$, $u,u'\in U$, is the link of $w_1\sigma$, and $\mathfrak g(uw_p\sigma u'^{-1})\ne0$ if $t\bar t\ne N^2$.
\end{enumerate}
\end{lemma}
\lean{New.Filling.NamesLink.frameProducts_top_top, New.Filling.NamesLink.isIntegralOff, New.Filling.NamesLink.link_eq, New.Filling.NamesLink.mem_entrySet, New.Filling.NamesLink.mem_floorGroup, New.Filling.NamesLink.reduction_eq, New.Filling.NamesLink.sq_lt_norm_sq}

\begin{proof}
(a) With $a$ the coordinate of $\sigma z$ at $z$ one has $t=h(z,\sigma z)=-Na$, and $-N=h(\sigma z,\sigma z)=\omega^*\omega-N\abs a^2$, since $\sigma z$ has no components along the hyperbolic planes; multiply by $N$. The difference
$t\bar t-N^2=N\omega^*\omega$ is nonnegative at each place, hence positive when it is not zero. (b) No $e^*_a$ has level below $\varrho$, and the Weyl element commutes with the change of basis entry by entry; the change of
basis on the right fixes $z$, so the column of $\sigma$ at $z$ remains, and in the lattice basis it is the vector $x-N^{-1}tz$, with $N=\pi\bar\pi$. (c) Each column of $w_p$ has one nonzero
entry; conjugating by the scales of the normalised basis multiplies the entries on the pair of floor $a$ by the squared scale or its inverse. (d) Both properties hold for $w_p\sigma$ and are invariant under $U$ on either
side (\cref{prop:links}(a)); the link of $w_p\sigma$ is that of $w_1\sigma$ by \cref{lem:link-isometry}(c), and $\mathfrak g(w_p\sigma)=\abs{d}^2(t\bar t-N^2)$, $d=\prod_ap_a$.
\end{proof}

\begin{proposition}[Class sums as orbit sums]\label{prop:orbit-sums}
Let the data be admissible, let the elements of $\mathcal T$ be units off a finite set $T$ of finite places, and let all terms over $(c,s,c')$ have nonzero leading minors
$D_1,\dots,D_\varrho$. Let $A$ be the block, for the cut function, of the terms over $(c,s,c')$ with $\mathfrak g\ne0$ and link $\Lambda$, and suppose $A\ne0$. Then there are $t\in\mathcal T$,
a rational $\sigma\in G^{(\varrho)}$ and $k\in K$ such that $t$ is the last diagonal entry of $Z(\sigma)$, $t\bar t-N^2$ is totally positive, $\Lambda$ is the link of $w_1\sigma$, $h_c^{-1}\sigma h_{c'}$
is integral off $T$, and $k$ reduces to $s$; and
\[
  A=\sum_{p\in\mathcal T^\varrho}\mathcal A^U\bigl(c,s,c';w_{p}\sigma\bigr),
\]
the sum over the heads $(p,t)$ of the orbit sums of the elements $w_p\sigma$.
\end{proposition}
\lean{New.Filling.exists_amplitudeSum_eq_sum_cosetAmplitude, New.Filling.translateByClassReps_inv_mul_mul, New.Filling.translateByClassReps_mem_floorGroup, New.Filling.orbitPoint_top_translateByClassReps_inv, New.Filling.mul_star_dotProduct_eq_last_frameProducts, New.Filling.translateByClassReps_weylElem_mul, New.Filling.leadingMinor_mem_boundedNonzeroIntegers, New.Filling.entry_mem_entrySet, New.Filling.isIntegralOff_of_mem_termsOver, New.Filling.isIntegralOff_weylElem_mul, New.Filling.sq_lt_norm_sq_last_frameProducts, New.Filling.injective_bigCellTerm, New.Filling.weylElem_padTuple, New.Filling.head_and_last_mem_entrySet, New.Filling.mem_termsOfLink_bigCell}
\lean{New.Filling.isStarStable_integralOff, New.Floors.IntegralPlace.modulus_evalPlaces_rpow_half}

\begin{proof}
\emph{The big cell.} A term with nonzero leading minors $D_1,\dots,D_\varrho$ lies in the big cell of the Bruhat decomposition for the flag $\langle e_1\rangle\subset\dots\subset A$: Gaussian elimination along the frame, with
the nonzero minors as pivots, writes $\gamma=uw_p\sigma u'^{-1}$ with $u,u'\in U(E)$ and $\sigma\in G^{(\varrho)}(E)$, and for fixed $\sigma$ the map $(p,u,u')\mapsto uw_p\sigma u'^{-1}$ is injective (the Weyl element depends only on the entries below
$\varrho$, so these may be completed by any last entry). The leading minor of order $l$ of $Z(\gamma)$ is that of $Z(w_p\sigma)$, the product $p_0\cdots p_{l-1}$ of the first entries, so $p_a=D_{a+1}/D_a$.

\emph{Entries.} The minors of a term with nonzero amplitude are nonzero integers with $\abs{D_l}_v^2\le b^{270n^2}$: at some $y$ both $W^\flat(y)$ and $W^\flat(\gamma^{-1}y)$ are nonzero, unipotent translates of the two
points lie in the boxes (\cref{prop:packing}(c)), the translates do not change the minors, and at points of the boxes the minors are bounded (\cref{lem:orbit-bound}(a)). Hence every entry $p_a$ lies in $\mathcal T$, and so
does $t=h(z,\sigma z)=D_{\varrho+1}/D_\varrho$.

\emph{The element $\sigma$.} Take a contributing term $\gamma_0=u_0w_{p_0}\sigma u_0'^{-1}$. Then $\mathfrak g(\gamma_0)=\abs{D_\varrho}^2(t\bar t-N^2)$, which is nonzero, so $t\bar t-N^2$ is nonzero and, by \cref{lem:weyl-norm}(a), totally
positive. The link of $\gamma_0$ is that of $w_1\sigma$ (\cref{lem:weyl-norm}(d)), and the link determines $\sigma$ up to the floor group's action by $U$ on both sides: by \cref{lem:link-isometry}(c), $w_p\sigma$ and $w_q\tau$ have the same link
exactly when $\sigma=\tau$. Hence every contributing term is $uw_p\sigma u'^{-1}$ with $p\in\mathcal T^\varrho$, each occurring once, and the block is the sum over the heads of the orbit sums. A lift $k\in K$ of $s$ exists because the
reduction is onto. The translate $h_c^{-1}\sigma h_{c'}$ is integral off $T$ because $h_c^{-1}\gamma_0h_{c'}$ is, and $h_c^{-1}u_0h_c$, $h_{c'}^{-1}u_0'h_{c'}$ and $w_{p_0}^{-1}$ are (the entries of $p_0$ and their inverses are units off $T$).
\end{proof}

\subsection{Kloosterman sums at a finite set of places}\label{ss:kloosterman}
Let $T$ be a finite set of finite places of $E$, stable under conjugation, and $E_T=\prod_{w\in T}E_w$, read as one ring with involution, integers $\OO_T$ and a Haar measure. Its unitary group
$G_T$ carries the congruence subgroup $K_T(M)$ of level $M$ in the group $K_T$ of integral points, the unipotent group $U_T$ and the character $\psi_{U}$. An element of the finite adelic
group is \emph{integral off $T$} if it and its inverse have integral entries outside $T$. The \emph{finite Bessel function} $\chi_M$ is $\psi_U(u)$ on $uK_T(M)$, $u\in U_T$, and $0$ off $U_TK_T(M)$.
For $x$, $s\in G_T$ the \emph{Kloosterman sum} is
\[
  \mathrm{Kl}(x,s)=\sum_{\substack{u,u'\in U_T/(U_T\cap K_T(M))\\ u^{-1}xu'\in sK_T(M)}}\chi_M(u)\overline{\chi_M(u')}.
\]
A \emph{filling datum at $T$} is a large lattice datum (\cref{ss:finite-data}) on $G_T$ whose modulus is the component of $M$ at $T$ and whose character is that of the data at $T$.
\lean{New.Filling.toPlaces, New.Filling.IsIntegralOff, New.Filling.IsFillingData, New.Floors.IntegralPlace.extendedBesselChar, New.Floors.PlaceData.UnipotentCosets, New.Floors.PlaceData.kloostermanSum, New.Filling.exists_isIntegralOff_toPlaces_eq, New.Filling.exists_isFillingData, New.Filling.exists_places_integralOff, New.Filling.eventually_mem_integralOff, New.Filling.exists_integralPlace, New.Filling.charOnPlaces_apply_eq_ψf, New.Filling.charOnPlaces_add_star}
\lean{New.Filling.ratToInf_weylElem, New.Filling.ratToInf_weylElem_mul, New.Filling.mem_unitaryLevel_iff, New.Filling.mem_latticeStabilizer_iff, New.Filling.exists_isAddHaarMeasure_skewAdjoint, New.Filling.exists_plancherel_kernel, New.Filling.forall_charOnPlaces_mul_eq_one_iff, New.Filling.locallyIntegrable_prod_one_add_sqrt, New.Filling.ofReal_genIndex_le_enorm_prod_sq, New.Filling.isMulLeftInvariant_haar_inv, New.Filling.star_neg_inv_deltaE}
\lean{New.Filling.IsFillingData.levelUnit_eq, New.Filling.IsFillingData.psi_eq, New.Filling.IsFillingData.toIsLargeLatticeData}

\begin{lemma}[The place $E_T$]\label{lem:place-T}
Let $T$ be finite and stable under conjugation.
\begin{enumerate}[label=\textup{(\alph*)}]
\item For finitely many nonzero elements of $E$ there is such a $T$ off which they are units and off which $M^{-1}$, $N^{-1}$ and the lattice basis with its inverse are integral; every finite adele is
integral off all sufficiently large $T$.
\item A finite adelic point is integral, respectively integral and congruent to $1$ modulo an idele $u$ with $u^{-1}$ integral off $T$, exactly when it is integral off $T$ and its component at $T$ is so.
\item The integers $\OO_T$ are self-dual for the character $x\mapsto\psi_{\rm f}(-x/\varpi)$ read at $T$, which is trivial on traces; there is a Haar measure on the skew elements of mass $1$ on the integral ones,
the unitary group $G_T$ is unimodular, and the generalised index of $x\in E_T$ is at most the square of $\prod_{w\in T}(1+(\norm{x_w}^{-1}-1)^{1/2})$.
\item For $x$ with $x$, $x^{-1}$ integral off $T$, the square root of the modulus of its component at $T$ is $\prod_v\abs x_v^{-1}$ (the product formula); in particular the modulus of the component of $M$ is $(M^f)^{-2}$.
\end{enumerate}
\end{lemma}
\lean{New.Filling.eventually_mem_integralOff, New.Filling.exists_places_integralOff, New.Filling.mem_unitaryLevel_iff, New.Filling.mem_latticeStabilizer_iff, New.Filling.exists_isAddHaarMeasure_skewAdjoint, New.Filling.forall_charOnPlaces_mul_eq_one_iff, New.Filling.charOnPlaces_add_star, New.Filling.isMulLeftInvariant_haar_inv, New.Filling.ofReal_genIndex_le_enorm_prod_sq, New.Floors.IntegralPlace.modulus_evalPlaces_rpow_half, New.Filling.IsFillingData.modulus_levelUnit}

\begin{proof}
(a) It suffices that $T$ contain the finitely many places where the given adeles are not integral, together with their conjugates. (b) Integrality is a condition at each finite place, and it splits into the
condition at $T$ and the one off $T$; off $T$ the congruence modulo $u$ follows from integrality, as $u^{-1}$ is integral there. (c) Self-duality reduces by extension by zero to that of $\hat\OO_E$ for $\psi_{\rm f}(-\cdot/\varpi)$,
$\varpi$ generating the different; the character is trivial on traces because $\varpi$ is skew and $\psi_{\rm f}$ takes the same value at conjugate adeles. With $\iota^2=-1$ and $\tau+\bar\tau=1$ integral at $T$, the map
$(a,b)\mapsto\tau\iota a+b$ from pairs of skew elements onto $E_T$ takes the product of a Haar measure on the skew elements with itself to the Haar measure, and the integral skew elements to the integral elements.
An orthonormal basis over $\hat\OO_E$ identifies the adelic group with that of the standard form, whose measure is bi-invariant; pushing it forward from the points integral off $T$, each point at $T$ having
such a lift, shows that the inverse of a left Haar measure is left invariant. The bound for the index is checked place by place. (d) The inequality $[x]^{1/2}\le\prod_v\abs x_v$ for $x$ and for $x^{-1}$, integral off $T$,
gives equality; apply it to $M$ squared.
\end{proof}

\begin{theorem}[Orbit sums are Kloosterman sums]\label{thm:kloosterman}
Let the data be admissible and $T$ a finite set of finite places, stable under conjugation, off which $M^{-1}$, $N^{-1}$ and the lattice basis are integral.
\begin{enumerate}[label=\textup{(\alph*)}]
\item The product $E_T$ is a finite place carrying a filling datum $F$ with $3\varrho+6$ further hyperbolic pairs and an orthonormal part of size $k-(6\varrho+13)$. Every point of $G_T$ is the
component at $T$ of a finite adelic point integral off $T$.
\item (Counting by measures.) Let $B$ be a subgroup of the finite adelic group and $\mu$ an s-finite measure on $G_T$, carried by the components at $T$ of the elements of $B$ integral off $T$
and right invariant under them. Then the order of the image of $B\cap K$ modulo $M$ times $\mu(K_T(M))$ is $\mu(K_T)$.
\item (Finite orbital integrals.) Let $Z$ be a group with $K_T(M)\le Z\le K_T$ and $W$ its \emph{Bessel vector}, with coordinates $W_a(y)=\chi_M(ya^{-1})$ at the cosets $a$ of $K_T(M)$ in $Z$. If
$ab^{-1}\in sK_T(M)$, then $\int_{G_T}W_a(y)\overline{W_b(x^{-1}y)}\,dy=\vol(K_T(M))\,\mathrm{Kl}(x,s)$.
\item (Orbit sums.) Let $\gamma=w_p\sigma$ with $\sigma\in G^{(\varrho)}(E)$, the inverses of the entries $p_a$ integral off $T$, and $h_c^{-1}\gamma h_{c'}$ integral off $T$. Let $x$ and $s$ be the components
at $T$ of $h_c^{-1}\gamma h_{c'}$ and of an integral point $k$ with reduction $\bar k$. Then
\[
  \mathcal A^U(c,\bar k,c';\gamma)=\mathrm{Kl}(x,s)\cdot\int_{G_\infty}W(y)\otimes W(\gamma^{-1}y)^*\,dy,
\]
with $W$ the full or the cut Bessel function at the infinite places.
\end{enumerate}
\end{theorem}
\lean{New.Filling.IsFillingData.card_reducedImage_mul_measure_congruenceSubgroup_eq, New.Floors.IntegralPlace.IsLargeLatticeData.besselOrbitalIntegral_besselVector_apply, New.Filling.IsFillingData.cosetAmplitude_eq_kloostermanSum_smul, New.Filling.IsFillingData.scaledOrbitalIntegral_smul_eq_cosetAmplitude, New.Filling.hasSum_besselOrbitalIntegral_cosetAmplitude, New.Filling.IsFillingData.mem_termsOver_iff, New.Filling.IsFillingData.exists_injective_mem_besselU_and_isIntegralOff, New.Filling.IsFillingData.extendedBesselChar_toPlaces_translateByClassReps, New.Filling.besselOrbitalIntegral_ratToInf_besselU_mul, New.Filling.IsFillingData.mem_principalLevel_iff, New.Filling.IsFillingData.reduction_eq_iff, New.Filling.IsFillingData.psi_simpleRootSum_eq_unipotentChar, New.Filling.IsFillingData.modulus_levelUnit, New.Filling.IsFillingData.unipotentOrder_mul_unipotentMass_eq_one, New.Floors.PlaceData.tsum_extendedBesselChar_mul_star_extendedBesselChar, New.Floors.PlaceData.inv_mul_mem_congruenceSubgroup_iff, New.Floors.PlaceData.unipotentCosets_eq_of_inv_mul_mem, New.Floors.PlaceData.exists_out_inv_mul_mem, New.Floors.IntegralPlace.extendedBesselChar_eq_of_inv_mul_mem, New.Filling.toPlaces_translateByClassReps_weylElem_mul, New.Filling.vectorOfLastEntry_evalPlaces}
\lean{New.integral_comp_quotientMk, New.inv_mul_out_coe_mem, New.Floors.PlaceData.inv_mul_out_mem, New.Floors.PlaceData.mem_unitaryLevel_one_of_inv_mul_mul_inv_mem}

\begin{proof}
(a) \emph{The datum.} In the lattice basis the Gram matrix of $X_0\perp Ez$ is integral with an integral inverse, over every
subring and so over $\OO_T$; hence the lattice of integral vectors at $T$ is unimodular. The Gram matrix of $X_0\perp Ez$ in the lattice basis splits at $T$ into $3\varrho+6$ further hyperbolic pairs and an orthonormal part of size $k-(6\varrho+13)$ (\cref{ss:finite-data}),
the subtraction being exact because $k$ is large. The rank conditions of a large datum hold because $k$ is large compared with $\varrho$ and with the degree of $E$, and the zeta value of $\OO_T$ is at most the global one
(\cref{ss:filling-data}); the character is that of \cref{lem:place-T}(c), and the modulus is the component of $M$. \emph{Lifting.} A point $y=1+z$ of $G_T$ has the lift $1+\phi(z)$, where $\phi$ extends each entry by zero;
the unitarity condition $z^*Jz+z^*J+Jz=0$ is preserved by $\phi$, and the lift is integral off $T$. If $y$ lies in $U_T$ or in $H_T$, so does the lift.
(b) Choose lifts $k_s\in B\cap K$ of the elements $s$ of the image of $B\cap K$ modulo $M$. Almost everywhere for $\mu$ the integral points at $T$ are the disjoint union of the cosets $K_T(M)k_s^{-1}$: a point of the support
of $\mu$ is the component of an element of $B$ integral off $T$, and if it is integral at $T$ it is integral (\cref{lem:place-T}(b)), hence lies in $B\cap K$ and in exactly one coset; by right invariance all the cosets have
measure $\mu(K_T(M))$.
(c) Put $yb$ for $y$; the integrand is then constant on the cosets $yK_T(M)$, because $\chi_M$ takes the same value at the points of a coset $yK_T(M)$ (it is $\psi(\sigma(u)/M)$ at $uk$ with $\sigma$ the sum of the entries at the
simple roots). For a left invariant measure and an open subgroup of finite volume, the integral of a function on the cosets is the volume of the subgroup times the sum over the cosets. In
$\sum_{y}\chi_M(ys^{-1})\overline{\chi_M(x^{-1}y)}$ a term is zero unless $yK_T(M)=usK_T(M)=xu'K_T(M)$ with $u$, $u'\in U_T$; these are unique modulo $U_T\cap K_T(M)$ (an integral $s$ normalises $K_T(M)$), and the value is
$\psi_U(u)\overline{\psi_U(u')}$. Hence $(u,u')\mapsto xu'K_T(M)$ is a bijection between the supports of the two sums that preserves the terms.
(d) \emph{Unfolding.} The group $U_c(M)\times U_{c'}(M)$ acts freely on the terms in $U\gamma U$, and each orbit contains one element $(r_w)^{-1}\gamma r'_{w'}$ for lists $r$, $r'$ of rational points of $U$, one in each coset $uU_c(M)$, $u'U_{c'}(M)$
for which $u^{-1}\gamma u'$ is a term; by \cref{prop:blocks}(b) the cutoffs may be replaced by the indicator of these representatives, so the orbit sum is the sum over the pairs of representatives giving a term of
the orbital integrals $\int W(y)\otimes W(\delta^{-1}y)^*dy$. \emph{Rational and local pairs.} By (a) and strong approximation, the translates $h_c^{-1}uh_c$ of the rational points $u$ of $U$ integral off $T$ match the cosets
$uU_c(M)$ with the cosets of $U_T\cap K_T(M)$ in $U_T$; in a term $u^{-1}\gamma u'$ over $(c,\bar k,c')$ the translates of $u$ and $u'$ are integral off $T$, and $(r_w)^{-1}\gamma r'_{w'}$ is a term over $(c,\bar k,c')$ exactly when
$w^{-1}xw'\in sK_T(M)$, by (b) of \cref{lem:place-T}. \emph{Phases.} For rational $u$, $u'$ of $U$, $J(u^{-1}\gamma u')=\psi_U(u)\overline{\psi_U(u')}J(\gamma)$ at infinity, the archimedean character of $U$ being, by the product formula for
$\psi$, the inverse of its finite part; and the finite Bessel function at the component at $T$ of the translate of $u$ is $\psi_U$ at the finite part of $u$, the character being trivial on the integral points off $T$ and
invariant under conjugation by $h_c$. Each pair therefore contributes $J(\gamma)$ times the term $\chi_M(w)\overline{\chi_M(w')}$ of the Kloosterman sum.
\end{proof}

\begin{corollary}[Orbit sums and the finite orbital integral]\label{cor:orbit-entry}
In the situation of \cref{thm:kloosterman}(d), let $K_T(M)\le Z\le K_T$, let $a$, $b$ be cosets of $K_T(M)$ in $Z$ with $k^{-1}ab^{-1}\in K_T(M)$, let $q$ and $l$ be the components at $T$ of $p$ and of $h_c^{-1}\sigma h_{c'}$, and let
$\kappa=\prod_{a<\varrho}\mathrm{mod}(q_a)^{e_a}$ be the head density. Then the entry $(a,b)$ of the scaled orbital integral of the Bessel vector of $Z$ at $w_ql^{-1}$, times $J(\gamma)$, is $\kappa\vol(K_T(M))$ times the
orbit sum $\mathcal A^U(c,\bar k,c';\gamma)$. Moreover the order of the reduction of $U(\hat\OO)$ modulo $M$ times the measure of $U_T\cap K_T(M)$ is $1$.
\end{corollary}
\lean{New.Filling.IsFillingData.unipotentOrder_mul_unipotentMass_eq_one, New.Filling.IsFillingData.scaledOrbitalIntegral_smul_eq_cosetAmplitude}

\begin{proof}
Both the entry at $T$ and the orbit sum contain the Kloosterman sum (\cref{thm:kloosterman}(c), (d)). For the last statement, by \cref{lem:place-T}(d) the floor $l<\varrho$ of $U_T\cap K_T(M)$ has volume $(M^f)^{-(2d_l+1)}$ with
$d_l=k+1+2(\varrho-l-1)$, and $\sum_{l<\varrho}(2d_l+1)=\varrho(2k+2\varrho+1)$, while the reduction of $U(\hat\OO)$ has $M^{f\varrho(2k+2\varrho+1)}$ elements (\cref{lem:reduced-groups}).
\end{proof}

\subsection{Siegel's formula as an upper bound}\label{ss:siegel-top}
For a hermitian matrix $\beta$ of size $k$ over $E$ and a function $f$ on the finite adelic vectors of rank $k$, the \emph{density} is $\delta(f,\beta)=\int_z\int_xf(x)\psi_{\rm f}(z(x^*x-\beta)/2)\,dx\,dz$, the
integral over the self-adjoint $z$. The \emph{Siegel constant} is
\[
  C_{\rm Sie}=\Bigl(\frac{\pi^k}{(k-1)!\,N^k}\Bigr)^f\abs{D_E}^{1/4}\vol(\hat\OO_E)^k,
\]
and the \emph{scaled Siegel constant} $C^\sharp_{\rm Sie}$ is $C_{\rm Sie}$ times the number of classes and the normalising factor of the local measures at $T$ (\cref{thm:siegel-top}(b)).
\lean{New.Filling.adelicFrequency, New.Filling.localFrequency, New.Filling.adelicFrequency_add, New.Filling.localFrequency_adelicFrequency_add, New.Filling.sub_adelicFrequency_mul_adeleOfLocal, New.Filling.ψf_adelicFrequency_mul, New.Filling.ψf_quadric_eq_mul, New.Filling.siegelConst, New.Filling.scaledSiegelConst, New.Filling.piNAt, New.Filling.piNAt_mul_star, New.Filling.x₀Part_mem_integralOff, New.Filling.orbitPoint_toPlaces_translateByClassReps}
\lean{New.Filling.continuous_adelicFrequency, New.Filling.continuous_localFrequency}

\begin{theorem}[Siegel's formula at a finite set of places]\label{thm:siegel-top}
Let the data be admissible, $T$ as in \cref{thm:kloosterman} and $F$ a filling datum at $T$.
\begin{enumerate}[label=\textup{(\alph*)}]
\item (Densities.) Let $f$ be $\Psi\ge0$ at $T$ times the indicator of the vectors integral off $T$, with $\Psi$ continuous of compact support. Then $\delta(f,\beta)$ is the product of an integral
of Gauss sums of the standard lattice over the $z$ whose part $-\varpi z$ at $T$ is integral, and of the density of $\Psi$ at $T$; and
\[
  \abs{\delta(f,\beta)}\le(1+b^{-6})\vol(\hat\OO_E)^k\int_{y^*y=\beta}\Psi(y)\,d\mu_\beta(y),
\]
the integral over the sphere at $T$ for its fibre measure.
\item There is $D>0$ with $C^\sharp_{\rm Sie}=D\cdot\#\{\text{classes}\}\cdot C_{\rm Sie}$ and, for every $Z$ with $K_T(M)\le Z\le K_T$, $\mathrm{mr}(Z)\le\mu_H(K_T(M))\,\mu(K_T(M))^2\,[Z:K_T(M)]\,D$, where
$\mathrm{mr}(Z)=\mu(Z)\mu(K_T(M))/\mathrm{fl}^2$ is the \emph{normalised mass} and $\mathrm{fl}$ the mass of $U_T\cap K_T(M)$.
\item (Siegel's formula as an upper bound.) Let $\Psi\ge0$ have compact support at $T$ and be invariant under small translations and under the elements of $K_T(M)$ that fix all basis vectors
outside $X_0$. Let $\abs t_v>N$ at every infinite place. For pairs $(c,v)$ of a class and a rational vector $v$ of $X_0$ with $N(v^*v)=t\bar t-N^2$ and $x=h_c^{-1}v$ integral off $T$,
\begin{multline*}
  \sum_{(c,v)}\Psi\bigl(x-N^{-1}tz\bigr)\le(1+b^{-6})\,\#\{\text{classes}\}\,C_{\rm Sie}\prod_v(\abs t_v^2-N^2)^{k-1}\\
  \times\abs{\pi_N}_T^{-1}\int_{y^*y=(t\bar t-N^2)/N}\Psi(y-N^{-1}tz)\,d\mu(y).
\end{multline*}
\end{enumerate}
\end{theorem}
\lean{New.Filling.normDensity_indicator_eq_mul, New.Filling.IsFillingData.enorm_normDensity_le_mul_lintegral_normFiberMeasure, New.Filling.IsFillingData.exists_scaledSiegelConst_eq_and_massRatio_le, New.Filling.siegel_constant_eq, New.Filling.IsFillingData.siegel_sum_le_lintegral_sphere, New.Filling.hasCompactSupport_testFunction, New.Filling.IsFillingData.testFunction_mulVec, New.Filling.skewUnit_inv_mul_lastEntrySphereIndex_piNAt, New.Filling.IsFillingData.classes_eq_and_eq_of_inv_mul_mem_congruenceSubgroup}

\begin{proof}
(a) \emph{Splitting the frequency.} A skew frequency $b$ at $T$ becomes a self-adjoint finite adele $\tilde b$ by extension by zero and multiplication by $-\varpi^{-1}$, and the map $z\mapsto-\varpi z$ at $T$ is a left
inverse of $b\mapsto\tilde b$, compatibly with translations; the difference $z-\tilde b$ for $b$ the part of $-\varpi z$ at $T$ vanishes at $T$. With $\tau+\bar\tau=1$ at $T$, $\psi_{\rm f}(\tilde by/2)=\psi(\tau by_T)$ for the character
$\psi$ of the place $E_T$. \emph{Splitting the vector.} Write $x=y+w$ with $y$ the extension by zero of a vector at $T$ and $w$ an adelic vector integral at $T$. The phase $\psi_{\rm f}(\zeta(x^*x-\beta)/2)$, for
$\zeta=\tilde b+z'$ with $-\varpi z'$ integral at $T$, is the product of a phase off $T$ (in $w$ and $z'$) and a phase at $T$ (in $y$ and $b$): the values of the form at $x$ and $w$ differ only at $T$, and the cross term is
integral at $T$, so its phase is $1$. Integrating, $\delta(f,\beta)$ is the product of the integral over $z'$ of the Gauss sums of the standard lattice and of the density of $\Psi$ at $T$. \emph{The bound.} The
Gauss sums of the standard lattice, integrated over the self-adjoint $z'$ with $-\varpi z'$ integral at $T$, are bounded by $\vol(\hat\OO_E)^k$ times the ideal zeta sum of $\hat\OO_E$ at $\mu=3\varrho+4$, since
$2(\mu+1)\le k$; this zeta sum is at most $1+2^{-\varrho}\le1+b^{-6}$, by the Euler product bound of the proof of \cref{lem:sublattice-zeta}(c), now at $\varrho+2\varrho+4$
with $[E\,{:}\,\Q]\le2^{2\varrho}$. At $T$ Fourier inversion in the skew frequency turns the density of $\Psi$ into the
integral over the sphere $y^*y=\beta$ for its fibre measure.

(b) In $\mathrm{mr}(Z)=\mu(Z)\mu(K_T(M))/\mathrm{fl}^2$ cover $Z$ by its $[Z:K_T(M)]$ cosets of $K_T(M)$, so $\mu(Z)\le[Z:K_T(M)]\mu(K_T(M))$, and use $o\cdot\mathrm{fl}=1$ for the order $o$ of the reduction of the unipotent group
(\cref{cor:orbit-entry}). Take $D=\#\bar R_0\,o=\#\bar H\,o^2$; the first identity is the definition of $C^\sharp_{\rm Sie}$. For the measure $\mu_H$ of $H$, the integral points have mass $1$, and by \cref{thm:kloosterman}(b) the
group $K_T(M)\cap H$ has index $\#\bar H$ in them, so $\#\bar H\,\mu_H(K_T(M))=1$; hence $\mathrm{fl}^{-2}=o^2=D\,\mu_H(K_T(M))$.

(c) \emph{The test function.} Let $f$ be the function on adelic vectors of $X_0$ that is $\Psi(x-N^{-1}tz)$ at $T$, in the lattice basis, times the indicator of the vectors integral off $T$. It has compact support (a vector
in the support is an integral vector plus the extension by zero of its part at $T$, read off the support of $\Psi$), and it is invariant under the level of the finite adelic unitary group of $X_0$, because at $T$ an
element of the level acts on the vectors with last entry $t$ as an element of $H\cap K_T(M)$. \emph{Siegel--Weil.} By the Siegel--Weil formula for the definite space $X_0$ (\cref{thm:B-siegel-formula}), the class sum of $f$ at
$\beta=(t\bar t-N^2)/N$ is the mass times the archimedean density $\prod_v\pi^k\beta_v^{k-1}/(k-1)!$ times $\delta(f,\beta)$. At the free level the stabilisers are trivial, so the weighted count of the class sum is the plain
count of the pairs $(c,v)$; each pair of the sum has norm $\beta$ and is integral off $T$, so it is a term of the class sum. The bound of (a), the evaluation of the class sum of $1$ (the number of classes), the product
formula for $N=\pi_N\bar\pi_N$, which gives the modulus of $\pi_N$ at $T$ as a power of $N^{-1}$, and the expansion of $C_{\rm Sie}$ and $\beta_v=(\abs t_v^2-N^2)/N$ give the claim.
\end{proof}

The \emph{fibre pairing} of two heads $\omega$, $\omega'$ at $T$ is the integral over $H\cap G^{(\varrho)}$, the fixer of all basis vectors outside $X_0$, of the product of the scaled orbital integrals of the
Bessel vector of $Z$ at $w_\omega l^{-1}$ and $w_{\omega'}l^{-1}$; it is a function of $e=l^{-1}z$, and its integral over the vectors $e$ with $h(z,e)=t$ is the \emph{density at the last entry}
$J_{\rm last}(\omega,\omega';t)$ (\cref{ss:top-floor}).
\lean{New.Floors.PlaceData.fiberPairing, New.Floors.PlaceData.lastEntryDensity, New.Floors.PlaceData.scaledOrbitalIntegral, New.Floors.PlaceData.headDensity}

\begin{lemma}[Classes over one point]\label{lem:classes-point}
Let $F$ be a filling datum at $T$ and $K_T(M)\le Z\le K_T$.
\begin{enumerate}[label=\textup{(\alph*)}]
\item Let $\gamma_1$, $\gamma_2$ be rational in the floor group with the same product $h(z,\gamma z)$ and the same part in $X_0$ of $\gamma z$, and let their left classes agree. If their translates by the class representatives
are integral off $T$ and lie in one coset of $K_T(M)$ at $T$, then $\gamma_1=\gamma_2$ and the right classes agree.
\item Let the heads $a_p$ have entries below $\varrho$ that are units off $T$, let $R$ be the Gram matrix of the tensors $\rho_p\int W(y)\otimes W((w_p\gamma)^{-1}y)^*dy$, and let $\rho_p$ times the head density be the same number $w$ for all
$p$. For integral $n_i$ in $Z$ at $T$ with distinct reductions $s_i$, the squared norms of $\sum_p\mathcal A^U(c,s_i,c';w_p\gamma)$, times $w^2$, $[Z:K_T(M)]$ and $\mu(K_T(M))^2$, add up to at most the sum against $R$ of the
pairings of the scaled orbital integrals at $T$ of two heads, at $g^{-1}$ for $g$ the component of $h_c^{-1}\gamma h_{c'}$.
\item If moreover $R$ is positive semidefinite, the classes $(c'_i,\bar n_i,\gamma_i)$ are distinct, the translates send $z$ to one point $e$ at $T$, and the left class, $h(z,\gamma_iz)$ and the part in $X_0$ of $\gamma_iz$
are the same for all $i$, then the squared norms of the class sums, times $w^2\mu_H(K_T(M))\mu(K_T(M))^2[Z:K_T(M)]$, add up to at most the real part of $\sum_{p,q}R_{pq}\,\mathrm{fp}(\omega_p,\omega_q;e)$.
\end{enumerate}
\end{lemma}
\lean{New.Filling.IsFillingData.classes_eq_and_eq_of_inv_mul_mem_congruenceSubgroup, New.Filling.IsFillingData.sum_sq_norm_sum_cosetAmplitude_le_headSumNormSq, New.Filling.IsFillingData.mul_sum_norm_sum_cosetAmplitude_sq_le_fiberPairing, New.Filling.orbitPoint_toPlaces_translateByClassReps, New.Filling.IsFillingData.exists_step_sum_fiberPairing}

\begin{proof}
(a) $\gamma_1^{-1}$ and $\gamma_2^{-1}$ give the same point $\gamma^{-1}z$, so $\gamma_1^{-1}\gamma_2$ fixes $z$ as well as the hyperbolic planes, hence lies in the floor group of level $\varrho+1$; its translate lies in the principal level,
which forces the right classes to agree and makes it unipotent with trivial part in the definite unitary group, hence trivial. (b) By \cref{cor:orbit-entry}, the entry $(b,A_ib)$ of the sum over the heads of the scaled
orbital integrals at $T$, with $A_ib$ the coset of $n_i^{-1}b$, is $w\mu(K_T(M))$ times the class sum at the reduction $s_i$; these entries are distinct as the $s_i$ are, and the sum against $R$ of the pairings of two heads is
the squared norm of the vector of all entries. (c) Group the classes by the pair of the right class and the rational element, and bound each group by (b); by (a) different pairs have positions in different cosets
of $K_T(M)$. Integrating over $H\cap G^{(\varrho)}$ and dividing by its measure of $K_T(M)$ turns the sum over the cosets into the fibre pairing at $e$.
\end{proof}

\begin{theorem}[Siegel's step]\label{thm:siegel-step}
Let the data be admissible, $T$ and $F$ as in \cref{thm:siegel-top}, and $Z$ a group with $K_T(M)\le Z\le K_T$. Let heads $a_p$, $p\in J$, have entries below $\varrho$ that are units off $T$
and last entry $t$ with $\abs t_v>N$ at every infinite place. Let $(c_i,\bar n_i,c'_i,\gamma_i)$, $i\in j$, be distinct, with $\gamma_i\in G^{(\varrho)}(E)$, $h(z,\gamma_iz)=t$, $h_{c_i}^{-1}\gamma_ih_{c'_i}$ integral off $T$ and
$n_i\in Z$ at $T$. Let $R$ be, for every $i$, the Gram matrix of the tensors $\rho_p\int W(y)\otimes W((w_p\gamma_i)^{-1}y)^*dy$, and let $\rho_p$ times the head density at $T$ be $r\sqrt G$ with
$G=\prod_v(\abs t_v^2-N^2)^{k-1}$. Then
\[
  r^2\sum_i\Bigl\|\sum_p\mathcal A^U(c_i,\bar n_i,c'_i;w_p\gamma_i)\Bigr\|^2\le(1+b^{-6})\,C^\sharp_{\rm Sie}\,\re\sum_{p,q}R_{pq}\,\frac{J_{\rm last}(a_p,a_q;t)}{\mathrm{mr}(Z)}.
\]
\end{theorem}
\lean{New.Filling.IsFillingData.sum_mul_norm_sum_cosetAmplitude_sq_le_lastEntryDensity, New.Filling.IsFillingData.mul_sum_norm_sum_cosetAmplitude_sq_le_fiberPairing, New.Filling.IsFillingData.sum_sq_norm_sum_cosetAmplitude_le_headSumNormSq, New.Filling.IsFillingData.exists_step_sum_fiberPairing}

\begin{proof}
Group the classes by the point $g_i$ of Siegel's sum, the left class and the part in $X_0$ of $\gamma_iz$. The translate of a term by the class representatives moves $z$ to the vector $x_i-N^{-1}tz$ with last entry $t$ over $g_i$
(\cref{lem:weyl-norm}(b)), which has norm $(t\bar t-N^2)/N$ and is integral off $T$. \emph{One point.} By \cref{prop:top-floor}(c) the positive part of the real part of $\sum_{p,q}R_{pq}\mathrm{fp}(\omega_p,\omega_q;\cdot)$ agrees on the top
quadric with a function $\Psi\ge0$ of compact support, invariant under small translations and under $H\cap K_T(M)$; by \cref{lem:classes-point}(c) its value at the point over $g_i$ bounds the classes there, times
$w^2\mu_H(K_T(M))\mu(K_T(M))^2[Z:K_T(M)]$ with $w=r\sqrt G$. \emph{Summing the points.} Siegel's formula (\cref{thm:siegel-top}(c)) bounds the sum of $\Psi$ over the points by $(1+b^{-6})n_XC_{\rm Sie}G\abs{\pi_N}_T^{-1}$ times the
integral of $\Psi$ over the sphere of $t$; almost every point of the sphere lies on the top quadric, where $\Psi$ is the real part of the sum of the fibre pairings, and by \cref{prop:top-floor}(a) the normalised integral over the
sphere is $\re\sum_{p,q}R_{pq}J_{\rm last}(a_p,a_q;t)$. \emph{Constants.} Dividing by $w^2\mu_H(K_T(M))\mu(K_T(M))^2[Z:K_T(M)]\ge w^2\mathrm{mr}(Z)/D$ (\cref{thm:siegel-top}(b)) and using $C^\sharp_{\rm Sie}=Dn_XC_{\rm Sie}$, the factor $G$ cancels against
$w^2=r^2G$.
\end{proof}

\subsection{Sieve systems}\label{ss:sieve-systems}
Let $C$ be a ring with a measure (a place, or a product of places), $u\colon C\to U$ a map to entries, $\psi$ a character of $C$ with $\psi(0)=1$, and $n\ge0$. A \emph{head} is a sequence of
entries. A \emph{sieve system} of length $n$ consists of forms $J_i(p,q)$, $0\le i\le n+1$, on pairs of heads, weights $\Lambda_i(p)\ge0$, a decay $\beta(\xi)\ge0$ for units $\xi$ and a mass bound $B_0$, such that:
\begin{enumerate}[label=\textup{(S\arabic*)}]
\item for $i\le n$, after giving both heads the entry $u(c)$ at position $i$, $c\mapsto\psi(\xi c)J_{i+1}$ is integrable; for $\xi=0$ its integral is $J_i(p,q)$, and for a unit $\xi$ its modulus is at
most $\beta(\xi)\Lambda_i(p)\Lambda_i(q)$;
\item for $i<n$, if the heads have entries $u(c)$, $u(c')$ at position $i$ and $\xi=\iota(c)-\iota(c')$ is a unit, then $\abs{J_{i+1}(p,q)}\le\beta(\xi)\Lambda_{i+1}(p)\Lambda_{i+1}(q)$;
\item the top form $J_{n+1}$ is continuous, and positive semidefinite on heads with a common entry at position $n$;
\item $\abs{J_0}\le B_0$.
\end{enumerate}
At a real place, $\beta(\xi)=\Theta\,b^{2000n^9}\ell_{k-n^7}(1/\abs\xi)$ with $\ell_j(r)=r^j\min(1,r)^6$, $\Theta$ the scaled volume ratio of the shape, and the weights are $\max(1,\abs\delta^{-1})^{n^6}$ over the
products $\delta=p_0\cdots p_l$, $l<j$, if these have modulus at most $b^{60n^2}$, and $0$ otherwise.
\lean{New.SieveSystem, New.dampedPower, New.headWeight, New.headWeight_nonneg, New.headAt, New.HasLocalBounds, New.Shape.scaledVolumeRatio}
\lean{New.SieveSystem.continuous, New.SieveSystem.integrable_and_integral_update, New.SieveSystem.norm_bot_le, New.SieveSystem.norm_le_of_isUnit_sub, New.SieveSystem.posSemidef, New.SieveSystem.psi_zero}

\begin{proposition}[Products, positivity, semicontinuity]\label{prop:sieve-products}
\begin{enumerate}[label=\textup{(\alph*)}]
\item Finitely many sieve systems on $\C$, indexed by a finite set $\tau$, and one sieve system on $C$ have a product, which is a sieve system on $\C^\tau\times C$. Its forms, characters,
weights and bounds are the products of those of the factors, each factor read on the components of the heads.
\item Sum the forms $J_{i+1}$ over all pairs of heads from one finite family, with a common entry at position $i\le n$. The sum is integrable in the common entry, its integral is the sum of the
$J_i$, and its twisted integrals are bounded by $\beta(\xi)$ and the squared sums of the weights.
\item Such a sum of top forms against a positive semidefinite matrix is nonnegative, and its real part is lower semicontinuous in a parameter at each point at which the heads are continuous.
\end{enumerate}
\end{proposition}
\lean{New.SieveSystem.pi, New.norm_prod_mul_le_prod_mul, New.posSemidef_prod_mul, New.SieveSystem.integrable_mul_update, New.SieveSystem.integrable_update, New.SieveSystem.integral_update, New.SieveSystem.norm_integral_mul_update_le, New.SieveSystem.integrable_sum_update, New.SieveSystem.integral_sum_update, New.SieveSystem.norm_integral_mul_sum_update_le, New.SieveSystem.integrable_and_integral_sum_update, New.SieveSystem.sum_nonneg, New.SieveSystem.lowerSemicontinuousAt_sum, New.SieveSystem.sum_nonneg_and_lowerSemicontinuousAt}

\begin{proof}
(a) Everything is checked factor by factor. (S1): for heads in the product, the integrand over a common entry $c=(c_v,c_C)$ is the product of the integrands of the factors, each integrable, and Fubini on the product measure
factors the integral, at $\xi=0$ into the product of the $J_i$ and at a unit $\xi$ into a product bounded by $\prod_v\beta_v(\xi_v)\Lambda_{i,v}\Lambda_{i,v}\cdot\beta_C(\xi_C)\Lambda_{i,C}\Lambda_{i,C}$, since $\abs{\prod_va_v\cdot a}\le\prod_v\abs{a_v}\abs a$. (S3): the
entrywise (Hadamard) product of positive semidefinite kernels is positive semidefinite (Schur's product theorem), and a product of continuous functions is continuous. (S2) and (S4) are products of the local bounds.
(b) Integrability and the integral follow from (S1) by linearity of the finite sum; for the twisted integral, $\abs{\sum_{p,q}\int\psi(\xi c)J_{i+1}}\le\beta(\xi)\sum_{p,q}\Lambda_i(p)\Lambda_i(q)=\beta(\xi)(\sum_p\Lambda_i(p))^2$.
(c) \emph{Positivity}, by descent on $i$: at $i=n$ the sum is the value of a positive semidefinite form at $(1,\dots,1)$, by (S3); below, by (S1) it is the integral over a common next entry of the sums for $i+1$, which are
nonnegative. \emph{Semicontinuity}, again by descent: at $i=n$ the sum is continuous in the parameter wherever the heads are, by (S3); below, it is the integral over a common next entry of nonnegative functions that are
lower semicontinuous in the parameter, and Fatou's lemma gives lower semicontinuity of the integral.
\end{proof}

\begin{proposition}[Finite sieve system]\label{prop:finite-sieve}
Let the data be admissible, $F$ a filling datum at $T$ with $M^{-1}$ integral off $T$, $S_g$ a set of chosen places above $p_0$ and $d_w$ nonzero vectors as in the definition of the fixer. There is a
group $Z$ with $K_T(M)\le Z\le K_T$ such that:
\begin{enumerate}[label=\textup{(\alph*)}]
\item $Z$ contains the component at $T$ of every integral element whose reduction lies in $\mathrm{fixer}(S_g,d)$;
\item the floor forms $J_j=\mathcal I_j/\mathrm{mr}(Z)$ of the Bessel vector of $Z$ for $j\le\varrho$ (\cref{ss:floor-forms}), followed by $J_{\varrho+1}=J_{\rm last}/\mathrm{mr}(Z)$ at the prime above $N$,
form a sieve system at $T$ for the character of the data, the decay $\xi\mapsto[\OO_T:\xi\OO_T]^{(k-n^7)/2}[\OO_T:\OO_T\cap\xi^{-1}\OO_T]^{-3}b^{n^9f}$, the weights with exponent $n^6$, and
$\abs{J_0}\le(1+b^{-6})b^{-6k\#S_g}$.
\end{enumerate}
\end{proposition}
\lean{New.Filling.IsFillingData.exists_intermediateGroup_sieveSystem, New.Filling.IsFillingData.exists_intermediateGroup_norm_floorForm_div_le, New.Filling.IsFillingData.massIntegral_full_le, New.Filling.norm_div_massRatio_le, New.Filling.IsFillingData.norm_floorForm_div_massRatio_le, New.Filling.IsFillingData.enorm_floorForm_zero_mul_le, New.Filling.IsFillingData.norm_floorForm_zero_div_massRatio_le, New.Filling.IsFillingData.posSemidef_lastEntryDensity, New.Filling.two_mul_M_pow_sq_pow_le, New.Filling.frequencyDecay, New.finHeadWeight}

\begin{proof}
Take $Z$ from \cref{prop:fixed-lines}, so that (a) holds and $\kappa(Z)\vol(K_T)b^{6k\#S_g}\le\kappa(K_T)\vol(Z)$ for the mass integral $\kappa(Z)=\int\abs{\int_Z\chi_M(yxy^{-1})dy}^2dx$. \emph{Descent.} For $j<\varrho$ floor $j$ is
the integral of floor $j+1$ over a common entry (\cref{thm:floor-descent}(b)), and at the top the integral of $J_{\rm last}$ over a common last entry is $\mathcal I_\varrho$ (\cref{prop:top-floor}(d)); the same integrals twisted by
$\psi(\xi c)$ are the floor forms with the weight $\psi(\xi Z_{jj})$. \emph{Decay.} By \cref{thm:finite-decay}, with $t=k-n^7$ and $K'=n^6$, both the twisted integrals and the forms at heads with entries whose frequency
$\xi=M^{-1}(p_i^{-1}-q_i^{-1})$ is a unit are at most $\mathrm{mr}(Z)$ times $[\OO_T:\xi\OO_T]^{(k-n^7)/2}[\OO_T:\OO_T\cap\xi^{-1}\OO_T]^{-6}(2/\abs M_T)^{n^6}$ times the finite head weights; the modulus of the level is $(M^f)^{-2}$
(\cref{lem:place-T}(d)), and $(2(M^f)^2)^{n^6}\le b^{n^9f}$ because $M\le b^{23}$; the decay keeps only the third power of the index of the denominator. \emph{Top.} $J_{\rm last}$ is continuous, and positive semidefinite on heads with a
common last entry, by \cref{prop:top-floor}(d). \emph{Mass.} $\mathcal I_0$ of the Bessel vector of $Z$, times $\mathrm{fl}^2$, is at most $\kappa(Z)$, both being integrals of $\abs{k_0}^2$ up to the factor $\mathrm{fl}^2$
(\cref{thm:floor-descent}(c)). The mass integral of the full group is at most $\vol(K_T)\vol(K_T(M))(1+b^{-6})$ by \cref{thm:finite-mass} with $Z=K_T$, $Q=2$ and $t=n$, since $\vol(K_T)\le1+2^{-n}$ and $2^{-n}+2^{-n}=2^{-\varrho}\le b^{-6}$.
Hence
\[
  \abs{J_0}\le\frac{\kappa(Z)}{\vol(Z)\vol(K_T(M))}\le\frac{\kappa(K_T)}{\vol(K_T)\vol(K_T(M))\,b^{6k\#S_g}}\le(1+b^{-6})\,b^{-6k\#S_g}.
\]
\end{proof}

\subsection{Poisson summation and the sieve}\label{ss:sieve}
For $\nu\in E$ write $[\nu]_v=\max(1,\abs\nu_v)^2$ at an infinite place and $[\nu]_{\rm f}$ for the norm of the denominator ideal of $\nu$. The \emph{constant of Poisson summation} in the entry $a$ is
$c_a=\vol(\hat\OO_E)/\prod_v\kappa_{v,a}^2$, with the squared frame scales $\kappa_{v,a}=\abs{\kappa_v}^{2(\varrho-a)}/N$ ($\kappa_v$ as in \cref{ss:filling-data}) at the infinite places.
\lean{New.denominator, New.mem_denominator, New.sum_absNorm_inv_sq_le, New.squareCorner, New.squareIndex, New.Filling.poissonConst, New.Filling.poissonConst_nonneg, New.Filling.frameScaleSq, New.Filling.ratPoint, New.Filling.charOnPlaces, New.Filling.globalChar, New.Filling.GlobalHead, New.Filling.unitOrOneAtPlaces, New.Filling.IsGlobalSieveSystem, New.Filling.BoundsDecay, New.Filling.scaledEntry, New.Filling.scaledHead, New.Filling.entryDensity, New.Filling.densitySum, New.Filling.decayConst, New.Filling.headProduct, New.Filling.globalWeight}

\begin{proposition}[Frequencies and one-sided Poisson summation]\label{prop:poisson}
Let $E$ be totally complex with $f$ infinite places.
\begin{enumerate}[label=\textup{(\alph*)}]
\item For every finite $X\subset E$, $\sum_{\nu\in X}\prod_v[\nu]_v^{-3}[\nu]_{\rm f}^{-3}\le(64\cdot38\,e^4)^f$.
\item For a finite set $X$ of nonzero $T$-integers and $\norm{c_v}^{-1}\le\abs{x_v}\le B$, $B\ge1$, the sum over $y\in X$ of the product of the decay factors at the infinite places and at $T$ is at most
$\prod_v\bigl(\abs{x_v}^k(bB^2)^3\bigr)$.
\item (One-sided Poisson summation.) Let $V$ be a set of $T$-integers, read at the infinite places scaled by $r_v>0$, and let $F\ge0$ be integrable and lower semicontinuous at these points, with Fourier transform at
each unit $\xi$ of modulus at most $\beta(\xi)M_0$. If $\varepsilon$ bounds the finite sums of $\beta$ over the nonzero $T$-integers, scaled by $2\pi i/(\varpi_vr_v)$, then
$\sum_{x\in V}F(x)\le\vol(\hat\OO_E)\prod_vr_v^{-2}\,(\varepsilon M_0+\int F)$.
\end{enumerate}
\end{proposition}
\lean{New.sum_height_inv_cube_le, New.Filling.sum_frequencies_le, New.Filling.frequency_term_le, New.Filling.sum_ratPoint_le_mul_add_integral_of_fourier_le, New.Filling.absNorm_denominator_le, New.Filling.measurePreserving_scale_evalPlaces, New.Filling.mem_ratIntegralOff_of_transform_ne_zero, New.Filling.extendToAdeles, New.Filling.integrable_extendToAdeles, New.Filling.lowerSemicontinuousAt_extendToAdeles, New.Filling.integral_stdAddChar_mul_extendToAdeles, New.exists_coset_label, New.sum_coset_weight_le, New.sum_weight_filter_denominator_eq_le}
\lean{New.denominator, New.mem_denominator, New.denominator_ne_bot, New.exists_mem_algebraMap_eq_mul, New.sum_absNorm_inv_sq_le, New.squareCorner, New.squareIndex, New.norm_lt_one_of_abs_lt, New.norm_sub_lt_one_of_squareIndex_eq, New.squareIndex_squareCorner, New.weight_le_weight_corner, New.pairwise_half_le_norm_squareCorner_sub, New.sum_weight_corner_le}

\begin{proof}
(a) \emph{Squares.} Cover $\C$ by the half-open squares of side $\frac12$ with corners $m/2$, $m\in\Z[i]$; two points of one square are at distance less than $1$, distinct corners at distance at least $\frac12$, and the weight
$\max(1,\abs w)^{-6}$ of a point is at most $64$ times that of the lower left corner of its square. The corners are $\frac12$-separated, so by packing their weights add up to at most $38$. \emph{Cosets.} If two elements of $X$
differ by an algebraic integer, a product over the places of squares of side $\frac12$ holds at most one of them, since the product of the absolute values of a nonzero algebraic integer is at least $1$; so the
archimedean weights $\prod_v\max(1,\abs\nu_v)^{-6}$ of such a set add up to at most $(64\cdot38)^f$. \emph{Denominators.} The numbers with a given denominator $\mathfrak b$ lie in $\mathrm N\mathfrak b$ cosets of $\OO_E$: with
$(a)=\mathfrak b\mathfrak c$, multiplication by $a$ maps them into $\mathfrak c$ and $\OO_E$ onto $(a)$, and $\mathfrak c/(a)$ has $\mathrm N\mathfrak b$ elements. Hence the sum is at most $(64\cdot38)^f\sum_{\mathfrak b}\mathrm N\mathfrak b\cdot\mathrm N\mathfrak b^{-3}$, and
$\sum_{\mathfrak b}\mathrm N\mathfrak b^{-2}\le\zeta_E(2)\le e^{2[E\,{:}\,\Q]}=e^{4f}$, by the Euler product with all coefficients $1$.
(b) Let $y\ne0$ be integral off $T$ and $r_v=\abs{c_vy}_v^{-1}$. The decay factor at $v$ is $r_v^k\min(1,r_v)^6$, and that at $T$ is $[\OO_T:y\OO_T]^{k/2}[\OO_T:\OO_T\cap y^{-1}\OO_T]^{-3}$. By the product formula, as $y$ is integral off $T$,
$[\OO_T:y\OO_T]\le\prod_v\abs y_v^2$, so $\prod_vr_v^k\,[\OO_T:y\OO_T]^{k/2}\le\prod_v\abs{c_v}^{-k}\le\prod_v\abs{x_v}^k$; and $\min(1,r_v)\le B\max(1,\abs y_v)^{-1}$, while the norm of the denominator of $y$ is at most
$[\OO_T:\OO_T\cap y^{-1}\OO_T]$. Hence the product of the decay factors is at most $\prod_v(\abs{x_v}^kB^6)$ times $\prod_v[y]_v^{-3}[y]_{\rm f}^{-3}$, and (a) with $64\cdot38\,e^4\le b^3$ gives the bound.
(c) Extend $F$ to the adeles: $G$ reads $F$ after scaling by $r$ at the infinite places and projecting to $T$ on the adeles integral off $T$, and is $0$ elsewhere. It is integrable, lower semicontinuous at the diagonal images
of the points of $V$, and invariant under the integral finite adeles that vanish at $T$. Scaling by $r$ at infinity and projecting to $T$ takes the measure on the adeles integral off $T$ to $\vol(\hat\OO_E)/\prod_vr_v^2$
times the product measure, the scaling having determinant $\prod_vr_v^2$. By Poisson summation for the lattice $E$ in $\Ad_E$, in its one-sided form for nonnegative lower semicontinuous integrable functions (approximate
$G$ from below by continuous functions of compact support), $\sum_{x\in E}G(x)\le\sum_{y\in E}\abs{\hat G(y/\varpi)}$, the dual lattice of $E$ being $\varpi^{-1}E$ for the standard character. By the invariance, $\hat G(y/\varpi)=0$ unless $y$ is a
$T$-integer (translation multiplies the transform by a value of the character, these values must be $1$, and $\hat\OO_E$ is its own dual); for a $T$-integer, $\abs{\hat G(y/\varpi)}=\vol(\hat\OO_E)\prod_vr_v^{-2}$ times the modulus of the
transform of $F$ at the point given by $y$ and the scales $2\pi i/(\varpi_vr_v)$. The term at $y=0$ is $\vol(\hat\OO_E)\prod_vr_v^{-2}\int F$, and the others are at most $\vol(\hat\OO_E)\prod_vr_v^{-2}\beta(\xi)M_0$, whose sum over the finitely many
relevant frequencies is at most $\varepsilon M_0$.
\end{proof}

For a head $p$ of length $i\le\varrho$ with entries in a finite set $A$ of units integral off $T$ together with their inverses, the \emph{global weight} is the product over the partial products $\delta=p_0\cdots p_l$, $l<i$, of
$\max(1,\norm{c_{l,v}\delta}_v^{-1})^{n^6}$ at each infinite place (or $0$ if $\norm{c_{l,v}\delta}_v>b^{60n^2}$) and of $[\OO_T:\delta\OO_T]^{n^6/2}$ at $T$ (or $0$ if $\delta\notin\OO_T$), where $c_{l,v}$ is the product of the squared scales of the
first $l+1$ entries.
\lean{New.Filling.headProduct_ne_zero, New.Filling.headProduct_mem_ratIntegralOff, New.Filling.headProduct_injective, New.Filling.prod_max_inv_norm_mul_sqrt_genIndex_le, New.Filling.inv_norm_prod_frameScaleSq_le, New.Filling.globalWeight_scaledHead, New.Filling.isIntegral_headProduct_and_norm_le, New.Filling.globalWeight_scaledHead_le, New.Filling.sum_globalWeight_le, New.Filling.one_le_pow_mul_prod_poissonConst, New.Filling.sq_sum_globalWeight_le, New.Filling.boundsDecay_prod, New.Filling.frameScaleSq_le_realM_pow, New.Filling.prod_abs_realM}

\begin{lemma}[Weights and decay]\label{lem:sieve-weights}
Let the data be admissible and $i\le\varrho$.
\begin{enumerate}[label=\textup{(\alph*)}]
\item If the global weight of $p$ is nonzero, the partial products are integers of $E$ of modulus at most $b^{120n^2}$ at every infinite place; the weight is at most $b^{120n^8fi}$, and a head is determined by its partial
products.
\item The weights of the heads of length $i$ add up to at most $b^{130n^9f}$, and the square of the sum is at most $b^{400n^9f}\prod_{a<i}c_a$.
\item $\prod_v\abs{\kappa_v}=(M\,\mathrm{rd}_E/(2\pi))^f$, and the sums of the decay over nonzero $T$-integers, at all scales $c_v/\kappa_{v,a}$ with $\norm{c_v}^{-1}\le\abs{\kappa_v}/M$, are at most $C_{\rm dec}\,b^{2002n^9f}$, with
$C_{\rm dec}=\prod_v\Theta\bigl(\abs{\kappa_v}^{2\varrho+1}/(MN)\bigr)^{k-n^7}$.
\end{enumerate}
\end{lemma}

\begin{proof}
(a) A nonzero weight forces $\delta\in\OO_T$ and, as $\delta$ is integral off $T$, $\delta$ integral; the archimedean condition and $\prod_{a\le l}\kappa_{v,a}^{-1}\le b^{60n^2}$ (from $N\le b^{60n}$ and $\abs{\kappa_v}\ge1$) give $\abs\delta_v\le b^{120n^2}$. By the
product formula, $\prod_v\max(1,\norm{c_v\delta}^{-1})\cdot[\OO_T:\delta\OO_T]^{1/2}\le R^f$ when $R$ bounds $\abs\delta_v$ and $\norm{c_v}^{-1}$; with $R=b^{120n^2}$ each partial product contributes at most $(b^{120n^2})^{fn^6}$, and there are $i\le n$ of
them. The entry $p_j$ is the quotient of two consecutive partial products. (b) By (a) and \cref{lem:box-integral}(a) the heads with nonzero weight are at most $((1+2b^{120n^2})^{2f})^{i}$ in number, and each weighs at most
$b^{120n^8fi}$; the product is at most $b^{130n^9f}$. The constants of Poisson summation satisfy $\prod_{a<i}c_a\ge b^{-100n^2f}$, by $\kappa_{v,a}\le b^{46\varrho}$ and $\vol(\hat\OO_E)=(2/\mathrm{rd}_E)^f$; so the square $b^{260n^9f}$ is at most
$b^{400n^9f}\prod_{a<i}c_a$. (c) The square of $\prod_v\abs{\varpi}_v$ is the norm of the different, $\abs{D_E}=\mathrm{rd}_E^{2f}$. For the sums apply \cref{prop:poisson}(b) with $x_v=\kappa_v^{2\varrho+1}/(MN)$, for which $\norm{c_v/\kappa_{v,a}}^{-1}\le\abs{x_v}\le B$
with $B$ a power of $b$ of exponent $O(n)$; then $b^{2000n^9f}$ from the archimedean decay, $b^{n^9f}$ from the finite one and $(bB^2)^{3f}$ add up to at most $b^{2002n^9f}$, the factors $\Theta\abs{x_v}^{k-n^7}$ making up $C_{\rm dec}$.
\end{proof}

\begin{theorem}[The sieve]\label{thm:sieve}
Let the data be admissible, $U\subseteq A$ finite sets of units, and $M^{-1}$ and $t$, $t^{-1}$ for $t\in A$ integral off $T$. Let the forms $J_{T,i}$ at $T$ form a sieve system of length $\varrho$ with mass
bound $m_T$, and let them be multiplied by the forms of the Bessel systems at the infinite places (\cref{def:bessel-system}), whose mass bounds are $(1+b^{-10})/\deg_v$. Then the real part of the sum of
the products of the top forms over all pairs of heads with entries in $A$ and a common last entry in $U$ is at most
\[
  \Bigl(\prod_{a\le\varrho}c_a\Bigr)\Bigl(\prod_v\frac{1+b^{-10}}{\deg_v}\,m_T+b^{2403n^9f}C_{\rm dec}\Bigr),
\]
where $C_{\rm dec}=\prod_v\Theta\bigl(\abs{\kappa_v}^{2\varrho+1}/(MN)\bigr)^{k-n^7}$ is the constant of \cref{lem:sieve-weights}(c).
\end{theorem}
\lean{New.Filling.sieve, New.Filling.densitySum_le, New.Filling.densitySum_succ_le_mul_densitySum_add, New.Filling.densitySum_succ_le, New.Filling.re_le_of_last_entries_ne, New.Filling.re_sum_sum_le_densitySum_add, New.Filling.entryDensity_scaledEntry, New.Filling.scaledHead_snoc, New.Filling.entryDensity_nonneg, New.Filling.lowerSemicontinuousAt_entryDensity, New.Filling.boundsDecay_prod, New.Filling.headProduct_ne_zero, New.Filling.headProduct_mem_ratIntegralOff, New.Filling.headProduct_injective, New.Filling.isIntegral_headProduct_and_norm_le, New.Filling.globalWeight_scaledHead_le, New.Filling.sum_globalWeight_le, New.Filling.one_le_pow_mul_prod_poissonConst, New.Filling.sq_sum_globalWeight_le, New.Filling.frameScaleSq_le_realM_pow, New.Filling.prod_abs_realM, New.Filling.decayConst_nonneg}
\lean{New.Filling.selfDualChar_apply, New.Filling.charOnPlaces_evalPlaces, New.Filling.continuous_charOnPlaces, New.Filling.continuous_globalChar, New.Filling.coe_stdAddChar_apply, New.Filling.wave_two_pi_I_mul, New.Filling.stdAddChar_div_deltaE, New.Filling.extendToAdeles_of_mem, New.Filling.extendToAdeles_of_notMem, New.Filling.extendToAdeles_toAdele, New.Filling.extendToAdeles_sub_adeleOfLocal_add, New.Filling.prod_max_inv_norm_mul_sqrt_genIndex_le, New.Filling.inv_norm_prod_frameScaleSq_le, New.Filling.globalWeight_nonneg, New.Filling.globalWeight_scaledHead}

\begin{proof}
By \cref{prop:sieve-products}(a) the products form a sieve system over the infinite places and $T$, for the character $e^{c-\bar c}$ at the infinite places times the character of $T$, with mass bound
$B_0=\prod_v(1+b^{-10})\deg_v^{-1}m_T$, the global weights and the product of the decays. Write $\Sigma_i$ for the \emph{density sums}: $\Sigma_0$ is the real part of the mass, and for $i>0$ $\Sigma_i$ is the real part of the
sum over $t\in A$ of the sum of $J_i$ over the pairs of heads of length $i$ with entries in $A$ and common last entry $t$, read at the scaled entry $t$ (scaled at $v$ by $\kappa_{v,i-1}$). The sum in the theorem is at most
$\Sigma_{\varrho+1}$, the terms being nonnegative (\cref{prop:sieve-products}(c)).

\emph{Removing one entry.} Let $i\le\varrho$. The sum of $J_{i+1}$ over the pairs of heads of length $i$ at a common next entry $x$ is a nonnegative integrable function of $x$, lower semicontinuous at the scaled
rational points (\cref{prop:sieve-products}(b), (c)); its integral is the sum of $J_i$ over these pairs, and its transform at a unit $\xi$ is at most $\beta(\xi)$ times the squared sum $\Gamma_i$ of the weights. One-sided Poisson summation
(\cref{prop:poisson}(c)) with the scales $r_v=\kappa_{v,i}$ gives
\[
  \Sigma_{i+1}\le c_i\Bigl(\varepsilon\Gamma_i+\re\sum_{p,q}J_i(p,q)\Bigr),
\]
with $\varepsilon=C_{\rm dec}b^{2002n^9f}$ bounding the decay over the frequencies (\cref{lem:sieve-weights}(c), with $c_v=2\pi i/\varpi_v$, for which $\norm{c_v}^{-1}\le\abs{\kappa_v}/M$). In the sum over all pairs of heads of length $i\ge1$, the
pairs with equal last entries give $\Sigma_i$, and those with different last entries $t$, $s$ have the frequency $(1/t-1/s)/M$, a nonzero $T$-integer, so by (S2) each is at most the decay there times the two weights; their sum is
at most $\varepsilon\Gamma_i$ (\cref{lem:sieve-weights}(c) with $c_v=M$). Hence $\Sigma_{i+1}\le c_i(\Sigma_i+2\varepsilon\Gamma_i)$.

\emph{Induction.} With $\Gamma=b^{400n^9f}$ we have $\Gamma_i\le\Gamma\prod_{a<i}c_a$ (\cref{lem:sieve-weights}(b)), and induction on $i$ gives
\[
  \Sigma_i\le\Bigl(\prod_{a<i}c_a\Bigr)\bigl(B_0+2i\Gamma\varepsilon\bigr),\qquad i\le\varrho+1.
\]
At $i=\varrho+1$, $2(\varrho+1)\Gamma\varepsilon\le b^{n^9f}b^{400n^9f}b^{2002n^9f}C_{\rm dec}=b^{2403n^9f}C_{\rm dec}$, since $2(\varrho+1)\le k\le b$.
\end{proof}

\subsection{The main term}\label{ss:main-term-bound}
The \emph{normalising constant} $r$ is the product over the infinite places of $\sqrt{c_G}$ and the powers of the squared frame scales, divided by the number of Bessel coordinates, where
$c_G=\vol\U(k)^2/(\vol\U(k+2\varrho+1)\vol\U(k-1))$ is the open cell constant.
\lean{New.Filling.normalizingConst, New.Filling.normalizingConst_sq_pos, New.Shape.openCellConst, New.Filling.gramScalar_mul_headDensity}

\begin{lemma}[The constants of the main term]\label{lem:main-constants}
Let the data be admissible, $k=k'+1$, $w=\vol(\hat\OO_E)^n(\#I/\Lambda)^2/c_G^f$ with $\Lambda$ the product of the powers of the squared frame scales, and $u=b^{-1}$.
\begin{enumerate}[label=\textup{(\alph*)}]
\item $\prod_{a\le\varrho}c_a=r^2w$, and for a head with entries units off $T$ and last entry $t$ with $\abs t_v>N$, the product of the scalars of the normalised orbital integrals at the infinite places and of the head
density at $T$ is $r\bigl(\prod_v(\abs t_v^2-N^2)\bigr)^{(k-1)/2}$.
\item $C^\sharp_{\rm Sie}\,w\le(1+u^9)\mathcal D_{\rm main}\deg\,\nu^2$ and $C^\sharp_{\rm Sie}\,w\,C_{\rm dec}\le b^{n^9f}\#I\,\mathcal D_{\rm id}\nu^2$.
\item If $\#S_g\le f$, then
\[
  (1+u^6)C^\sharp_{\rm Sie}\prod_ac_a\Bigl(\prod_v\frac{1+u^{10}}{\deg_v}\cdot\frac{1+u^6}{b^{6k\#S_g}}+b^{2403n^9f}C_{\rm dec}\Bigr)\le r^2(1+u^4)\,b^{-6k\#S_g}\mathcal D_{\rm main}\nu^2.
\]
\end{enumerate}
\end{lemma}
\lean{New.Filling.one_add_inv_pow_ten_pow_le, New.Filling.one_add_pow_six_mul_add_pow_ten_le_one_add_pow_four}

\begin{proof}
(a) One factor $\kappa_{v,a}$ splits off each power in $\Lambda$, and the rest is the definition of $r$; for the head, the product formula cancels the entries of the head between the infinite places and $T$, and splitting off the
scale of the last frame vector in $r$ leaves $(\prod_v(\abs t_v^2-N^2))^{(k-1)/2}$. (b) The first is the bound for the scaled Siegel constant required by admissibility (\cref{prop:filling-admissible}), rewritten. For the second,
write $\mathrm{rd}_E=s^2$, $Q=Ms^2/(2\pi)$, so that $\prod_v\abs{\kappa_v}=Q^f$ and $\vol(\hat\OO_E)=(2/s^2)^f$. Apart from a common factor both sides are $f$th powers, of $A$ and of $b^{n^9}B$, and at one place $A$ is $B$ times
$(2/s^2)^{n+2\varrho^2}(s/s^{2\varrho})x^{-n^7}$ with $x=Q^{2\varrho+1}/(MN)$; these factors are at most $b^{n^9}$, as $2/s^2\le b$, $s/s^{2\varrho}\le1$ and $x^{-1}\le b^{23}b^{60n}$.
(c) The main part is at most $(1+u^9)^2(1+u^6)r^2b^{-6k\#S_g}\mathcal D_{\rm main}\nu^2$ by (a), (b) and $(1+u^{10})^f\le1+u^9$ (as $f\le b/1000$). The error part is at most $b^{2403n^9f}b^{n^9f}b^{2kf}\mathcal D_{\rm id}\nu^2r^2$, as $\#I\le b^{2kf}$;
since $k\ge512000n^9$ and $\#S_g\le f$, $2403n^9f+n^9f+2kf+10+6k\#S_g\le11kf$, and $b^{11kf}\mathcal D_{\rm id}\le\mathcal D_{\rm main}$ bounds the error by $u^{10}r^2b^{-6k\#S_g}\mathcal D_{\rm main}\nu^2$. Finally
$(1+u^6)\bigl((1+u^9)^2(1+u^6)+u^{10}\bigr)\le1+u^4$ for $0\le u\le10^{-3}$, all the powers on the left being at most $u^4/100$.
\end{proof}

\begin{theorem}[Main term]\label{thm:main-term}
Let $E$ carry Weil data, the filling data be admissible, $S_g$ a set of chosen places above $p_0$ and $d_w$ nonzero vectors.
\begin{enumerate}[label=\textup{(\alph*)}]
\item (Band limits.) Let $\sigma\in G^{(\varrho)}(E)$ have last frame product $t$. The product over the infinite places of the top forms of the Bessel systems at the heads $(p,t)$, $(q,t)$ is the
Gram matrix of the orbital integrals $\int W(y)\otimes W((w_p\sigma)^{-1}y)^*dy$ of the cut Bessel function, each weighted by
\[
  \prod_v\sqrt{c_G}\,(\abs{x_\varrho}^2-1)^{(k-1)/2}\abs I^{-1}\prod_{a<\varrho}\abs{x_a}^{k+1+2(\varrho-a-1)},
\]
with $x$ the scaled entries of the head at $v$.
\item For finitely many distinct $(c_i,\bar n_i,c'_i,\sigma_i)$ with $\sigma_i\in G^{(\varrho)}(E)$, last entries $t_i$ in a finite set of units off $T$ with $\abs{t_i}_v>N$, positions integral off $T$ and reductions in
$\mathrm{fixer}(S_g,d)$,
\[
  \sum_i\Bigl\|\sum_{p\in\mathcal T^\varrho}\mathcal A^U(c_i,\bar n_i,c'_i;w_p\sigma_i)\Bigr\|^2\le(1+b^{-4})\,b^{-6k\#S_g}\,\mathcal D_{\rm main}\,\nu^2.
\]
\item The link second moment of the terms with $\mathfrak g\ne0$ over the open elements of $\mathrm{fixer}(S_g,d)$ is at most $(1+b^{-4})b^{-6k\#S_g}\mathcal D_{\rm main}\nu^2$.
\end{enumerate}
\end{theorem}
\lean{New.Filling.linkSecondMoment_isOpenElement_le, New.Filling.sum_norm_sum_cosetAmplitude_sq_le, New.Filling.IsFillingData.mul_sum_norm_sum_cosetAmplitude_sq_le, New.Filling.prod_bandLimit_eq_gram, New.Filling.sieve_constant_le, New.Filling.scaledSiegelConst_mul_decayConst_le, New.Filling.prod_poissonConst_eq, New.Filling.one_add_inv_pow_ten_pow_le, New.Filling.one_add_pow_six_mul_add_pow_ten_le_one_add_pow_four}

\begin{proof}
(a) At each infinite place both heads end in the last frame product of $\sigma$, so by \cref{thm:local-band-limits}(b) the top form is the inner product of the normalised orbital integrals; the scaled orbital integrals of the
product vector are the products of the local ones, and the product over $v$ of the inner products is the Gram matrix of the products.

(b) \emph{Places.} Enlarge $T$ so that it contains the chosen places, the places above $N$ and $M$ and those where the lattice basis or the entries of the heads are not integral, and is stable under conjugation
(\cref{lem:place-T}(a)); take a filling datum at $T$ (\cref{thm:kloosterman}(a)) and $Z$ from \cref{prop:finite-sieve}, which contains the components of the $n_i$ because their reductions lie in the fixer. \emph{Siegel's step.} Group the
classes by their last entry $t$. For each $t$, \cref{thm:siegel-step} applies with $R$ the product of the top forms at the infinite places, which by (a) is the Gram matrix of the weighted orbital integrals, and with
$\rho_p$ the product of the scalars, for which $\rho_p$ times the head density is $r\sqrt G$ by \cref{lem:main-constants}(a). It bounds $r^2$ times the sum by $(1+b^{-6})C^\sharp_{\rm Sie}$ times the real part of the sum over the
heads of the product of the top forms at the infinite places and of $J_{\rm last}/\mathrm{mr}(Z)$ at $T$. \emph{The sieve.} Summed over $t$ this is the sum of \cref{thm:sieve} for the system of \cref{prop:finite-sieve}, with
$m_T=(1+b^{-6})b^{-6k\#S_g}$, so $r^2$ times the left side is at most the left side of \cref{lem:main-constants}(c). Dividing by $r^2>0$ gives the bound.

(c) The left side is a sum of squared norms of blocks; it suffices to bound finite sums of nonzero blocks. Let $s$ be open in the fixer. By \cref{lem:big-cell} every term over $(c,s,c')$ lies in the big cell, so by
\cref{prop:orbit-sums} each nonzero block of one link is a sum of orbit sums over the heads with entries in $\mathcal T$, for one $\sigma$. A block has two classes, a reduction and a link; distinct blocks give distinct tuples
$(c,c',\bar n,\sigma)$, as $\bar n$ is the reduction and $\sigma$ determines the link. Choosing $T$ so that the elements of $\mathcal T$ are units off $T$ and $M^{-1}$, $N^{-1}$ are integral off $T$, (b) applies to the finite family of
nonzero blocks, with last entries in $\mathcal T$.
\end{proof}
